\subsection{Modulus measure of a differential form}

Notation. --- We denote \(\mu\) a Haar measure (INT, VII, §1, no. 2) on the additive group of K. We denote \(\mu^{\otimes n}\) the measure \(\mu \otimes \cdots \otimes \mu\) onto \(K^n\) and \(\mu_{\mathbf{U}}^{\otimes n}\) its restriction to an open set U of \(K^n\) . For \(a \in K\) we denote by mod(a) the modulus of a (INT, VII, §1, no. 10 and AC, VI, §9, no. 1). If K = R, we have mod(a) = \textbar a\textbar; if K = C, we have mod(a) = \textbar a\textbar\^{}2; if K is ultrametric, we have mod(a) = \(q^{-v(a)}\) , where \(q\) denotes the number of elements in the residue field of K \(v\) the normalized valuation of K (AC, VI, §9, no. 1, prop. 1).

10.1.1. Let U and V be open subsets of \(K^{n}\) and let f: \(U \to V\) a map of class \(C^{r}\) , with \(r \in N_{K}\) . Let \(x \in U\) ; we call the Jacobian of f at x, and denote by \(\operatorname{Jac}_{x}(f)\) , the determinant of the linear map \(Df(x)\) derivative of f at x. We denote \(\operatorname{Jac}(f)\) the function \(x \mapsto \operatorname{Jac}_{x}(f)\) the function \(\operatorname{mod}(\operatorname{Jac}(f))\) is a continuous function in U with positive real values.

10.1.2 ("Change of variable in integrals"). Under the hypotheses of 10.1.1, suppose that \(f\) let be a class isomorphism \(\mathbf{C}^{r}\) of the manifold U onto the manifold V. Then the image under \(f\) of the measure mod(Jac(f)). \(\mu_{\mathrm{U}}^{\otimes n}\) is \(\mu_{\mathrm{V}}^{\otimes n}\) ; for every continuous function with compact support \(\varphi: \mathrm{V} \to \mathbf{C}\) , one has

\[
\int_ {\mathrm{V}} \varphi (y) \mu^ {\otimes n} (y) = \int_ {\mathrm{U}} \varphi (f (x)). \mathrm{mod} (\operatorname{Jac} _ {x} (f)) \mu^ {\otimes n} (x).\tag{1}
\]

10.1.3. Let X be a variety of class \(\mathbf{C}^r\) and let A be a subset of X. One says that A is locally negligible if, for every chart \((\mathrm{U},\varphi ,\mathrm{K}^n)\) of the manifold X, the set \(\varphi (\mathrm{A}\cap \mathrm{U})\) is \(\mu^{\otimes n}\)-negligible (INT, IV, § 2, no. 2); this condition does not depend on the choice of \(\mu\) , and it suffices to verify it for a family of charts whose domains cover A. Every set contained in the union of a countable family of locally negligible sets is locally negligible.

Examples

\begin{enumerate}
\def\labelenumi{\alph{enumi})}
\item
  Every submanifold of X that has codimension \(\geqslant1\) at every point (i.e. has empty interior) is locally negligible.
\item
  Let g be a function of class \(C^{r}\) on X, and let \(A_{g}\) the set of points \(x \in X\) such that \(g(x) = 0\) ; suppose that the interior of \(A_{g}\) is empty, and that \(r = \omega\) ; then \(A_{g}\) is locally negligible.
\item
  Let \(f: Y \to X\) be a morphism of manifolds of class \(C^{r}\) , the manifold \(Y\) being a countable union of compact sets. It is assumed that one has \(\dim_{y} Y \leqslant \dim_{f(y)} X\) for all \(y \in Y\) . Its image under \(f\) of every locally negligible subset of \(Y\) is a locally negligible subset of \(X\) . The same is true of the image under \(f\) of the set of points where \(f\) is not étale (“first Sard theorem”); in particular, if \(\dim_{y} Y < \dim_{f(y)} X\) for all \(y \in Y\) , \(f(Y)\) is a locally negligible subset of \(X\) .
\item
  Suppose K has characteristic zero. Let \(f\colon Y\to X\) be a morphism of manifolds of class \(\mathbf{C}^{r}\) ( \(r\in N_{K}, r\geqslant\infty\) ), Y being a countable union of compact sets. Let C be the set of points of Y where \(f\) is not a submersion (“critical points”). Then \(f(\mathbf{C})\) is a locally negligible subset of X (“second Sard theorem”). \(^{(1)}\)
\end{enumerate}

10.1.4. Let X be a paracompact manifold of class C'. On X there exists a class M of equivalent measures (INT, V, § 5, n° 6), and only one, such that, for every \(\nu \in \mathbf{M}\) and for every chart \(c = (\mathrm{U}, \varphi, \mathrm{K}^n)\) of X, the image under \(\varphi\) of the restriction of \(\nu\) to U is equivalent to \(\mu_{\varphi(\mathrm{U})}^{\otimes n}\) . One says that M is the canonical class of measures on X. For a subset A of X to be locally negligible (10.1.3), it is necessary and sufficient that it be locally negligible (INT, IV, § 2, n° 2) for one (resp. every) measure \(\nu \in \mathbf{M}\) .

If K = R and \(r \neq \omega\) (resp. K = R or C and \(r = \omega\) , resp. K is ultrametric), there exists \(\nu \in M\) such that, for every chart \(c = (\mathrm{U}, \varphi, \mathrm{K}^{n})\) of X, the density of \(\varphi(\nu_{\mathrm{U}})\) with respect to \(\mu_{\varphi(\mathrm{U})}^{\otimes n}\) is everywhere \textgreater0 and of class \(C^{r-1}\) (resp. of class \(C^{\infty}\) , resp. locally constant). If \(\nu\) and \(\nu'\) are two such measures, the density of \(\nu'\) with respect to \(\nu\) is everywhere \textgreater0 and of class \(C^{r-1}\) (resp. of class \(C^{\infty}\) , resp. locally constant).

10.1.5. Let X be a manifold of class \(\mathbf{C}^r\) ; let \(\mathrm{T}'(\mathrm{X})\) the cotangent bundle of X (8.2.2) and set \(\Omega = \det(\mathrm{T}'(\mathrm{X}))\) (7.9.9). The vector bundle \(\Omega\) has rank 1 at each point; when X is pure of dimension \(n\) , it is identified with the vector bundle \(\operatorname{Alt}^n(\mathrm{T}(\mathrm{X}); \mathrm{K}_{\mathbf{x}})\) . Let \(\omega\) a section of \(\Omega\) on X; one says, by abuse of language, that \(\omega\) is a differential form of maximum degree on X. Let \(c = (\mathrm{U}, \varphi, \mathrm{K}^n)\) be a chart of X; denote by \(u^1, \ldots, u^n\) the coordinate functions on \(\mathrm{K}^n\) . There exists a function \(f_c\) and only one on \(\varphi(\mathrm{U})\) such that \(\omega|_{\mathrm{U}} = \varphi^*(f_c.du^1 \wedge \cdots \wedge du^n)\) . We say that \(\omega\) is locally of integrable modulus if, for every chart \(c = (\mathrm{U}, \varphi, \mathrm{K}^n)\) of X, the real function mod \((f_c)\) : \(\varphi(\mathrm{U}) \to \mathbf{R}\) is locally integrable for \(\mu_{\varphi(\mathrm{U})}^{\otimes n}\) (INT, IV, § 4, n° 1); it suffices to verify this property for the charts of an atlas of X. If \(\omega\) is continuous, and in particular if \(\omega\) is of class \(\mathbf{C}^s\), with \(s \in \mathbb{N}_{\mathbf{K}}, s \leqslant r - 1\), then \(\omega\) is locally of integrable modulus.

10.1.6 (“Positive measure defined by a differential form of maximum degree”). Keep the notations of 10.1.5 and suppose moreover that X is separated and that \(\omega\) is locally of integrable modulus. If \(c = (\mathrm{U}, \varphi, \mathrm{K}^n)\) is a chart of X, denote by \(\nu_c\) the measure on \(\varphi(\mathrm{U})\) product of the measure \(\mu_{\varphi(\mathrm{U})}^{\otimes n}\) by the function mod( \(f_c\) ) (INT, V, § 5, n° 2, def. 2) and let \(\alpha_c\) the image of \(\nu_c\) by \(\varphi^{-1}\) . There exists on X a measure \(\alpha\) and only one such that, for every chart \(c = (\mathrm{U}, \varphi, \mathrm{K}^n)\) of X, the restriction of \(\alpha\) to U is equal to \(\alpha_c\) . One says that the measure \(\alpha\) is the modulus of \(\omega\) and one denotes it mod( \(\omega\) ) \(_\mu\) . It is a positive measure. When K = R and \(\mu\) is Lebesgue measure, one writes \(|\omega|\) instead of mod( \(\omega\) ) \(_\mu\) .

If X is purely of dimension n, and if a is a real number \textgreater0, one has \(\text{mod}(\omega)_{a\mu} = a^{n} \text{ mod}(\omega)_{\mu}\) .

Let A be a subset of X. For A to be locally negligible (10.1.3), it is necessary and sufficient that, for every open set U of X and every differential form \(\omega\) of maximum degree on U which is locally of integrable modulus, the set A ∩ U be locally negligible for the measure mod( \(\omega\) ) \(_{\mu}\) .

Suppose moreover X is paracompact and let \(\nu\) be a measure on X belonging to the canonical equivalence class M (10.1.4). The measure mod \((\omega)_{\mu}\) is based on \(\nu\) (INT, V, § 5, no. 2, def. 2). For mod \((\omega)_{\mu}\) to belong to M, it is necessary and sufficient that the set of \(x \in \mathbf{X}\) such that \(\omega(x) = 0\) be locally negligible.

10.1.7. Example. --- Let G be a Lie group over K, of finite dimension n, and let \(\omega\) be a differential form of degree n on G, invariant under left translations, and nonzero. The corresponding measure \(\text{mod}(\omega)_{\mu}\) is then a left Haar measure on the locally compact group G.

When G is the multiplicative group \(K^{*}\) , one may take for \(\omega\) the form dx/x, and one has \(\text{mod}(\omega)_{\mu} = (\text{mod})^{-1}.\mu\) , cf.~INT, VII, § 1, no. 10, prop. 14.
\begin{enumerate}
\def\labelenumi{(\arabic{enumi})}
\setcounter{enumi}{6}
\tightlist
\item
\end{enumerate}
% semantic-fragments: legacy-input-seam
\relax{}
